% !TeX root = ../papers/nc-p3.tex
\section{単一固有値の全Jordan窓とBell直線の完結}
\label{sec:split-single-eigenvalue-jordan}
\begin{theorem}\label{thm:split-single-eigenvalue-jordan}
平方固有値が一つの直線型可逆分解窓のうち、非半単純なものの無印添字保持同型類は次で尽きる。
全て標準Koszul左右AS正則、最小周期1であり、十五次元成分には属さない。
\begin{center}
\begin{tabular}{cl}
鎖分割 & 同型類\\\hline
$4$ & 一点\\
$3+1$ & $b\in k$, $b\sim-b$\\
$2+2$ & $(u,v)\in k\times k^\times$, $(u,v)\sim(-u,v^{-1})$、および別の一点\\
$2+1+1$ & $a\in k^\times$, $a\sim a^{-1}$、別の一点、さらに $p\in k^\times$, $p\sim-p$
\end{tabular}
\end{center}
各行と各行内の別欄は互いに交わらない。全て量子閉軌道へ流入する。
$2+1+1$ の $a$ 欄は前節の一辺の表を $s=1$ に延ばしたものであり、他の量子閉軌道との
照合にも同じ辺と三角積による判定が使える。
\end{theorem}

\begin{proof}[斉一な移送と四つの中心化群]
四次元模型を二つの標準二文字空間の積とし、形式を $B=J\otimes J$ とする。
通常の四テンソルの第一切断行列を右から $B$ 倍して作用素として見る。
左右の局所射影群は二つの $\operatorname{PSO}(B)$、$\tau=(02)(13)$ は随伴を取る。
共通倍率と左特殊直交作用により $P$ の作用素を
$L=(1+N)^{1/2}=1+N/2-N^2/8+N^3/16$ に取れる。ここで $N$ は自己随伴、$N^4=0$ である。
これは定理\ref{thm:split-jordan-transport}の極分解で、単一固有値1と $\det P=1$ に規格化したものである。
$1+N$ の固有値は1なので、$P$ の射影安定群 $H$ は
$N$ を中心化する特殊直交群を $\{\pm1\}$ で割ったものの左右対角埋込みである。
実際、作用素の射影倍率の二乗が1となり、左右Gramの等式から右側は $N$ と可換、
左側も射影的に同じとなる。

双対を各二文字で $J$ により同定する。二逆テンソルの作用素は、それぞれ $L^{-1}$ と
$L^{-1}(1-N^2/2)=(1-N)^{1/2}$ である。長さ $m$ の鎖を $k[z]/z^m$、
$N$ を $z$ 倍、対合を積の $z^{m-1}$ 係数として
\[
 U_t f(z)=(1+tz)^{(m-2)/2} f\big(z/(1+tz)\big)\pmod {z^m}
\]
と置く。$U_tU_h=U_{t+h}$ であり、$U_t$ は直交ユニポテントで
$U_tNU_t^{-1}=N/(1+tN)$ を満たす。直交性は、積の最上係数を留数
$\operatorname{Res} f(z)g(z)\,dz/z^m$ と書き変数変換すれば従う。
よって $U_1(1-N)^{1/2}U_1^{-1}=L^{-1}$ である。
第一の四対応は双対同定後に恒等、第二の四対応は二位置反復で
$g=U_1h$ の二つの射影二文字因子を取ればよい。全ての解はこれと $h\in H$ で得られる。
下記の通り $U_1$ は $H$ と可換なので、印付き同型は $h$ の通常の共役である。

長さ4の一鎖なら、それと可換な行列は $N$ の多項式、直交条件はその二乗が1であるため $H=1$。
$3+1$ では三鎖に $1,z,z^2$、一鎖に $f$ を取り $\langle f,f\rangle=1$ とする。
$H$ は $1\mapsto f$, $f\mapsto-z^2$, $z,z^2\mapsto0$ なる $D$ の
$\exp(bD)$ 全体である。実際、$1$ の像を $c+c_1z+c_2z^2+b f$、
$f$ の像を $e z^2+d f$ とすると、直交性と行列式1より
$c=d=\pm1$, $c_1=0$, $2cc_2=-b^2$, $ce=-bd$ を得る。全符号を除けば主張通りである。
$U_t$ は定数項を保ち、$z^2,f$ を固定するので $D$ と可換である。

以下 $N^2=0$ では $U_t=1$ である。$E$ を二次の上三角行列単位 $E_{12}$ とする。
$2+2$ は $N=E\otimes Z$ とその二因子交換、$Z=\operatorname{diag}(1,-1)$ である。
両者は像の極大全等方平面の二つの族に対応する。$\operatorname{PSO}(B)=\mathrm{PGL}_2^2$ で
$E,Z$ を逆の倍率で保つ必要があり、後者の倍率は $\pm1$ に限る。
従って $H=(\mathbb G_a\times\mathbb G_m)\rtimes C_2$ で、対合は $(u,v)\mapsto(-u,v^{-1})$ である。
非恒等成分は平方写像の全射性から単一共役類となる。
$2+1+1$ は $N=E\otimes E$ である。二つの上三角群の対角倍率が互いに逆になり、
$H=\mathbb G_a^2\rtimes\mathbb G_m$、$t$ の重みは $1,-1$ である。
$t\ne1$ では二つの平行移動を共役で除ける。$t=1$ なら零ベクトル、二つの非零座標軸、
積 $p=uv\ne0$ ごとの一類に分かれる。これで印付きの全解を得た。
\end{proof}

\begin{proof}[一移動と分解の交換]
$S$ を二つの二文字因子を交換する行列とし $K=L-1$ と書く。
分解交換後の第一交代成分は $S L^{-1}S^{-1}$ である。実際 $\pi_0=(03)(12)$ は
第一切断の両側の交換と各側での因子交換であり、双対同定後の第一逆行列は自己随伴である。
一移動と分解交換を同時に行ったときは $S(1-K)S^{-1}$ となる。
実際、自己随伴作用素の空間はスカラーと $\mathfrak{sl}_2\otimes\mathfrak{sl}_2$ の直和で、
四文字の巡回は因子交換と第一因子だけの交代形式による随伴を取る。後者は跡零行列に負号を付ける。
$K$ は跡零自己随伴なので表示式を得る（体積の共通符号は射影的に消える）。

どちらの操作でも第二の四対応は $Sg^{-1}S^{-1}$ の二位置反復となる。
実際、双対同定後の四つの行列を $d_i$ とすると、分解交換では $d_{\pi_1(i)}^{-1}$、
一移動と交換では旧第二対応を新第一の座標に用いて $d_{i-1}^{-1}$ になる。
今は $d_{i+2}=d_i$ なので両者は同じである。
さらに $U_{1/2}(1-K)U_{1/2}^{-1}=L^{-1}$ である。これは $N^4=0$ の上で
$(1+N/(1+N/2))^{1/2}-1=K/(1+K)$ を展開した等式である。
$U_{1/2}$ は $g$ と可換だから、これを $S$ で運んで左右同じ基底変更に使えば、
二窓が一致する。従って無印最小周期は全て1である。
\path{computations/verify_split_single_eigenvalue_jordan.py} は各鎖の反対角形式、両移送、
通常二進座標での巡回・交換および $3+1$ の中心化作用を直接検算する。

指定分解の一径数を含む安定群の階数は、その中心化群の階数であり、$1+\operatorname{rank}C_H(h)$ である。
六ブロックの相対倍率の核が共通法線倍率で、残りは第一の四対応と $h$ の比較から従う。
階数1なら原始余指標は共役を除き符号だけである。$2+2$ の恒等成分内の $h$ では階数2だが、
第二因子の対角トーラスは奇位置の能動二文字と偶位置の相補二文字という一つの色にだけ非自明に作用する。
他の色の二文字は同じ重みなので、二文字ずつの二重みを与える余指標は元の分解とその逆だけである。

交換は $N$ を $S[-N/(1+N)]S^{-1}$ に運ぶ。$-N/(1+N)$ は元と同じ核の旗を持ち、
$\det S=-1$ なので、全鎖長が偶数のとき極大全等方平面の二族が交換される。
従って長さ4の二向きから一点、$2+2$ では片方の向きの印付き分類をそのまま得る。
$3+1$ では $1,z,z^2,f$ に対する $D_0=\operatorname{diag}(1,-1,1,1)$ を使う。
$D_0 U_{-1}S^{-1}$ は特殊直交で新しい $N$ を元に戻す。$D_0$ は $D$ を固定し $U_t$ を
$U_{-t}$ に変えるので、移送 $g$ の新しいパラメータは $-b$ である。
$2+1+1$ の非零平行移動では階数1で、交換は $(u,v)$ をまず $(-v,-u)$ にする。
さらに $-E\otimes E$ を $E\otimes E$ に戻す二対角倍率の積は $-1$ なので、
二座標軸は一類、積は $p\sim-p$ となる。
\end{proof}

\begin{proof}[残りのトーラスと他の行との比較]
$2+1+1$ の対角な $h$（恒等を含む）では、
通常二進座標で $P=P(1)+E_{0011}$ として $L_i^1=T(a,1/a)_i F_i$ を使える。
実際、この $P$ の第一切断は単位行列に一つの上三角行列単位を加えた形である。
$B$ を一側に作用させて随伴作用素に戻せば $1+N/2$ であり、二つの因子のビット変更と
対角倍率で $N=E\otimes E$ にできる。この変更と双対の $J$ はいずれも単項なので、
前節の二逆転置に $s=1$ を代入して得る上の解は、$H$ の中でも単項、従って上三角の中の対角になる。
$a$ の倍率を変えると両因子の対角比は互いに逆の任意の倍率で変わるので、対角な $h$ を尽くす。
全位置で相異なる四指標を持つ極大トーラスがあるから、前節の一辺の無印同型証明もそのまま使える。
この辺は $02$、全三角積は $-a$ である。置換 $(02)$ が全三角積を逆数にするので $a\sim a^{-1}$ となる。

階数1あるいは一つの色にだけ追加トーラスが作用する行では、分解は上で述べた通り一意であり、
$P,Q$ の鎖分割は等しい。相互の異なる鎖分割や半単純な $P$ と混同されない。
四指標が相異なる一辺の行では前節の表とも同じ有限置換判定で比較できる。
具体的に、向き付き四面値を $T_{123},1/T_{023},T_{013},1/T_{012}$ と並べると、
$s=1$ では互いに逆な二値が二個ずつである。前節の一辺でこの形を持つには既に等しい
第二・第四値以外に第一・第三値も等しく、従って $s^4=1$ が必要である。
$s^2=-1$ なら二値の積は $-1$ であって条件を満たさない。置換による全値の同時反転もこの区別を保つ。

最後にJordan鎖の標準重みで $N$ を零にすれば、$4,3+1$ では $g$ の次数付きは恒等、
$2+2$ では第一因子の平行移動が消え第二因子のトーラス元または法線化群の反転が残る。
$2+1+1$ では両因子で $E$ の重みが正なので平行移動は消える。
従って残る二因子はいずれも半単純である。$L=1$ でのこの二因子は、
通常座標 $P(1)$ に戻すと定理\ref{thm:split-spectral-line-classification}の $CJ,DJ$ の共役類である。
実際、第一二位置の $J$ で $B$ を単位行列にし、相補基底で第一の四対応を恒等に戻すと
最初の二行列は各因子の左に $J^{-1}$ を掛けたものとなる。
よって閉軌道は単一平方固有値の量子類であり、定理\ref{thm:split-line-component-incidence}により
十五次元成分とは交わらない。
\end{proof}
